\section{Introduction}
\label{sec:introduction}
Knowledge Tracing (KT) estimates a learner's evolving mastery state from sequential interactions and predicts future correctness \citep{corbett1994knowledge,piech2015deep,abdelrahman2023knowledge}. Recent KT models, including memory-based, attention-based, and transformer-style variants, are typically compared using aggregate predictive metrics such as the area under the ROC curve (AUC) and accuracy (ACC) \citep{zhang2017dynamic,ghosh2020context,liu2023simplekt}. These metrics are useful, but they can hide systematic differences on low-frequency knowledge components (KCs), especially when most test interactions are concentrated on dense, frequently observed KCs.

This issue is not merely a reporting detail. In educational decision support, predicted probabilities can be used to decide whether a learner should receive remediation, practice another item, or progress to a new concept. A model that ranks answers well overall may still be poorly calibrated on sparse or limited cold-start KCs. In such cases, high predicted confidence can be misleading even when the model's aggregate AUC appears competitive. Calibration metrics such as Expected Calibration Error (ECE), Brier score, and reliability diagrams therefore provide a complementary diagnostic view \citep{guo2017calibration,brier1950verification,degroot1983comparison}.

This paper treats KT evaluation itself as a reproducibility and diagnostic problem. We do not propose a new KT architecture, nor is this protocol designed to establish a new state-of-the-art model or to change the ``winning'' model choice. Its purpose is to diagnose reliability properties---in particular calibration on sparse concepts---that aggregate discrimination metrics may conceal. We specify a dataset-agnostic diagnostic protocol that extends standard KT evaluation with train-only KC-frequency stratification, per-stratum calibration diagnostics, concept-level cold-start analysis, and an operational leakage and predictive-sanity workflow (L1--L8). The framing is intentionally conservative: the goal is reusable evaluation infrastructure for future KT studies---including graph-based and self-supervised approaches whose gains are often claimed on sparse concepts---rather than any claim of state-of-the-art performance. The protocol is not only preventative. In our own instantiation, its predictive-sanity channel flagged a prediction--label misalignment in the temporal evaluation pipeline that aggregate metrics alone would not have revealed (Section~\ref{sec:discussion}); we report the audit trail as part of the study.

The scientific question is not whether sparse concepts universally perform worse. Under the conditions of this study, sparse training evidence does not universally degrade discrimination, but it can expose dataset-dependent calibration vulnerability. We characterize the conditions associated with that vulnerability and provide reproducible diagnostics for identifying when it matters. The instantiation supports six bounded findings: (1)~lower KC frequency does not imply universal AUC degradation; (2)~calibration can be more sensitive than aggregate discrimination in some sparse-concept regimes; (3)~the calibration--frequency relationship is dataset-dependent; (4)~observable factors such as difficulty, item support, learner exposure, and curriculum position explain part of this heterogeneity; a controlled reduction of training rows for the same originally dense KCs does not universally increase ECE, though SimpleKT on Junyi Academy shows a large within-KC increase; (5)~sparse-concept diagnostics are therefore conditionally important rather than equally informative on every KT log; (6)~L1--L8 is a reproducible audit workflow, and L8 has empirical value in detecting prediction--label misalignment and the controlled fault classes reported in Appendix~\ref{app:l8_fault}.

\subsection{Motivating Diagnostic Gap}
Under standard benchmarking protocols, overall KT performance can be misleading because sparse KCs and limited cold-start KCs can exhibit different predictive and calibration behavior from the dense majority. For example, a model may achieve competitive overall AUC on a learner-based split while showing higher calibration error, or unstable AUC estimates, on sparse KC strata---or, on another dataset, no such sparse-stratum degradation at all. This discrepancy matters when predicted probabilities are used to make decisions regarding student remediation, progression, or mastery thresholds. If a model's probabilities are poorly calibrated on the concepts where evidence is scarcest, educational interventions may be triggered inappropriately, despite a deceptively high population-level discrimination score. The gap is therefore diagnostic rather than architectural: evaluation should report performance and reliability at the KC-frequency level, with occupancy and sample-size context, and with leakage audits that keep pipeline artifacts from being read as scientific findings.

In addressing this gap, this paper makes three contributions:
\begin{itemize}
    \item \textbf{Problem 1 (Hidden Stratum Behavior):} Overall AUC/ACC values can mask heterogeneous predictive behavior or metric instability on low-frequency concepts.
    \textbf{Contribution 1:} We formalize a train-only KC-frequency stratification protocol---with an explicit strict cold-start group ($f_{\mathrm{train}}(c) = 0$) separated from very sparse concepts---that prevents sparse-bucket leakage and enables sample-size-aware KT diagnostics across pre-registered frequency buckets.
    \item \textbf{Problem 2 (Unreliable Probability Estimates):} Aggregate rankings do not reveal whether predicted probabilities correspond to empirical correctness rates, which is crucial for educational decision-making.
    \textbf{Contribution 2:} We evaluate calibration-oriented metrics---ECE, the Brier decomposition, and reliability diagrams---across frequency buckets, and we report when discrimination and calibration diverge. When predicted probabilities are treated as a global mastery/remediation gate, we further simulate stratum-wise false-mastery and miss rates at a locked threshold (Section~\ref{sec:threshold_simulation}). The associated empirical claim is bounded: calibration can be more sensitive than aggregate discrimination in some sparse-concept regimes, the calibration--frequency relationship is dataset-dependent, and a population threshold can then produce a larger false-mastery rate on sparse KCs than on dense KCs for SimpleKT on ASSISTments 2012 (positive $\Delta\mathrm{FM}$ on all five training seeds, four unique partitions; Limited occupancy).
    \item \textbf{Problem 3 (Hidden Evaluation Pipeline Contamination):} KT evaluation pipelines can suffer from subtle, undocumented data leakage and alignment errors across folds or splits.
    \textbf{Contribution 3:} We operationalize eight named leakage and predictive-sanity checks into a reproducible KT evaluation workflow with verifiable artifacts. Channels L1--L7 instantiate standard evaluation safeguards (split hygiene, train-only definitions, validation-only selection) as inspectable checks; they are not offered as a new auditing methodology. Channel L8 is a warm-cohort predictive-sanity test that, in this study, flagged a temporal prediction--label misalignment (Appendix~\ref{app:alignment}). A subsequent controlled fault-injection validation (Appendix~\ref{app:l8_fault}) shows that the same signals are sensitive to some tested alignment-fault classes (label-identity shifts; global row-index mismatch) and insensitive to others (adjacent-timestep prediction shifts that remain autocorrelated). We do not claim that L8 detects all evaluation bugs.
\end{itemize}

Our experimental investigation is structured around four Research Questions (RQs):
\begin{itemize}
    \item \textbf{RQ1 (Stratum-level Predictive Heterogeneity):} How does predictive performance vary across KC-frequency strata, and is lower training frequency associated with systematic AUC degradation across datasets?
    \item \textbf{RQ2 (Calibration Profiles by Frequency):} How do model calibration profiles change across different concept frequency strata, and is calibration more sensitive than aggregate discrimination in some sparse-concept regimes?
    \item \textbf{RQ3 (Explanatory Conditions):} Under what dataset and concept conditions is lower training frequency associated with meaningful calibration degradation, and when are sparse-concept diagnostics informative rather than uniformly required?
    \item \textbf{RQ4 (Limited Cold-start Concept Diagnostics):} How do baseline KT models behave on KCs with zero or limited training-fold frequency under temporal splits?
\end{itemize}

RQ1 exercises Contribution 1, RQ2 exercises Contribution 2 (including the single-fold GKT/CL4KT instantiation and simulated gate), RQ3 uses Contributions~1--2 together with covariate-adjusted associations and a controlled sparsification experiment (Appendix~\ref{app:a9_sparsify}), and RQ4 draws jointly on Contributions 1 and 3. The remainder of this paper is organized as follows. Section~\ref{sec:related} reviews KT evaluation, sparse-concept analysis, and calibration. Section~\ref{sec:protocol} presents the diagnostic protocol. Section~\ref{sec:experiments} reports the experimental design and diagnostic results. Section~\ref{sec:discussion} discusses implications and limitations.
